% =====================================================================
\chapter{Integrales dobles y la integral gaussiana}
\label{tema:6}
% =====================================================================

La integral de una variable suma los valores de una función a lo largo de un intervalo; la integral doble los suma sobre una región del plano y mide, por ejemplo, el volumen bajo una superficie. Su cálculo se reduce a dos integrales de una variable sucesivas (teorema de Fubini), y en regiones circulares conviene usar coordenadas polares. La aplicación más importante para el curso es el cálculo de la integral gaussiana $\int_{-\infty}^{\infty}e^{-x^2}dx=\sqrt\pi$, que no puede obtenerse con primitivas elementales y sí con un rodeo por el plano. La presentación sigue las sesiones~47--54 de 18.02SC~\cite{mit1802}, las §5.1--5.3 de OpenStax~\cite{openstax3}, las §11.1--11.14 de Apostol~\cite{apostol2} y las §14.1--14.2 de Strang~\cite{strangcalc}.

\section{Integral doble sobre un rectángulo}

Sea $R=[a,b]\times[c,d]$ y $f$ acotada en $R$. Se divide $[a,b]$ en $m$ subintervalos de longitud $\Delta x$ y $[c,d]$ en $n$ de longitud $\Delta y$, lo que divide $R$ en $mn$ rectángulos $R_{ij}$ de área $\Delta A=\Delta x\,\Delta y$, y se elige un punto $(x_{ij}^*,y_{ij}^*)$ en cada uno.

\begin{definicion}[Integral doble]\label{def:integral-doble}
La \emph{integral doble} de $f$ sobre $R$ es
\[
\iint_R f(x,y)\dA=\lim_{m,n\to\infty}\sum_{i=1}^{m}\sum_{j=1}^{n}f(x_{ij}^*,y_{ij}^*)\,\Delta A ,
\]
si el límite existe y no depende de la elección de los puntos; se dice entonces que $f$ es \emph{integrable} en $R$. Toda función continua en $R$ es integrable.
\end{definicion}

Si $f\ge0$, cada sumando es el volumen de un prisma de base $R_{ij}$ y altura $f(x_{ij}^*,y_{ij}^*)$, y la integral es el \emph{volumen} del sólido comprendido entre $R$ y la gráfica de $f$. Si $f=1$, la integral es el área de $R$.

\begin{proposicion}[Propiedades]\label{prop:propiedades-integral}
Para $f,g$ integrables en $R$ y $\alpha,\beta\in\R$:
\begin{enumerate}
  \item linealidad: $\iint_R(\alpha f+\beta g)\dA=\alpha\iint_Rf\dA+\beta\iint_Rg\dA$;
  \item monotonía: si $f\le g$ en $R$, entonces $\iint_Rf\dA\le\iint_Rg\dA$;
  \item aditividad: si $R$ se divide en dos regiones $R_1$, $R_2$ que solo comparten frontera, $\iint_R f=\iint_{R_1}f+\iint_{R_2}f$.
\end{enumerate}
\end{proposicion}

Las tres propiedades se heredan de las sumas de Riemann al pasar al límite: por ejemplo, si $f\le g$, cada suma de $f$ es menor o igual que la correspondiente de $g$. El cálculo efectivo se apoya en el teorema siguiente, cuya demostración puede verse en Apostol~\cite[§11.10--11.11]{apostol2}.

\begin{teorema}[Fubini, para rectángulos]\label{teo:fubini}
Si $f$ es continua en $R=[a,b]\times[c,d]$, entonces
\[
\iint_R f\dA=\int_a^b\left(\int_c^d f(x,y)\,dy\right)dx=\int_c^d\left(\int_a^b f(x,y)\,dx\right)dy .
\]
\end{teorema}

En la integral interior $\int_c^df(x,y)\,dy$, la variable $x$ se trata como constante. Geométricamente, es el área de la sección del sólido por el plano vertical de abscisa $x$, y el volumen se obtiene sumando secciones.

\begin{corolario}[Integrandos separables]\label{cor:separable}
Si $f(x,y)=g(x)h(y)$ con $g$ y $h$ continuas, entonces $\iint_{[a,b]\times[c,d]}g(x)h(y)\dA=\left(\int_a^bg\right)\left(\int_c^dh\right)$.
\end{corolario}

\addtocounter{definicion}{1}%

\section{Regiones generales y cambio del orden de integración}

\begin{definicion}[Regiones elementales]\label{def:regiones}
Una región es \emph{de tipo~I} si es de la forma $\{a\le x\le b,\ g_1(x)\le y\le g_2(x)\}$, y \emph{de tipo~II} si es $\{c\le y\le d,\ h_1(y)\le x\le h_2(y)\}$, con $g_i$, $h_i$ continuas.
\end{definicion}

Para $f$ continua, el teorema de Fubini se extiende a estas regiones~\cite[§11.12--11.14]{apostol2}:
\begin{align*}
\iint_Df\dA&=\int_a^b\int_{g_1(x)}^{g_2(x)}f(x,y)\,dy\,dx&&\text{(tipo I)},\\
\iint_Df\dA&=\int_c^d\int_{h_1(y)}^{h_2(y)}f(x,y)\,dx\,dy&&\text{(tipo II)}.
\end{align*}
Los límites de la integral interior pueden depender de la variable exterior; los de la exterior han de ser constantes. El método práctico es: dibujar la región, fijar la variable exterior y leer, sobre un segmento vertical (tipo~I) u horizontal (tipo~II), dónde entra y dónde sale de la región.

\addtocounter{definicion}{1}%

Cuando la integral interior no tiene primitiva elemental, puede bastar con cambiar el orden. Hay que describir la \emph{misma} región del otro modo.

\addtocounter{definicion}{1}%

\section{Coordenadas polares}

En polares, $x=r\cos\theta$, $y=r\sin\theta$, con $r\ge0$. Las regiones naturales son los \emph{rectángulos polares} $\{r_1\le r\le r_2,\ \theta_1\le\theta\le\theta_2\}$, sectores de corona circular. Su área es la diferencia de dos sectores circulares:
\[
\tfrac12r_2^2\Delta\theta-\tfrac12r_1^2\Delta\theta=\tfrac{r_1+r_2}{2}\,(r_2-r_1)\,\Delta\theta=\bar r\,\Delta r\,\Delta\theta ,
\]
donde $\bar r$ es el radio medio. Así, un rectángulo polar pequeño tiene área aproximada $r\,\Delta r\,\Delta\theta$, no $\Delta r\,\Delta\theta$: los rectángulos más alejados del origen son más grandes. Sumando y pasando al límite se obtiene el resultado siguiente~\cite[§11.26--11.31]{apostol2}.

\begin{teorema}[Cambio a coordenadas polares]\label{teo:polares}
Si $f$ es continua en la región $D=\{\alpha\le\theta\le\beta,\ r_1(\theta)\le r\le r_2(\theta)\}$, con $0\le\beta-\alpha\le2\pi$, entonces
\[
\iint_Df(x,y)\dA=\int_\alpha^\beta\int_{r_1(\theta)}^{r_2(\theta)}f(r\cos\theta,r\sin\theta)\,r\,dr\,d\theta .
\]
\end{teorema}

El factor $r$ es el valor absoluto del determinante de las derivadas parciales de $(x,y)$ respecto de $(r,\theta)$, el \emph{jacobiano} $\bigl|\begin{smallmatrix}\cos\theta&-r\sin\theta\\\sin\theta&r\cos\theta\end{smallmatrix}\bigr|=r$, de acuerdo con la interpretación del determinante como factor de cambio de área (\cref{teo:det-area}).

\addtocounter{definicion}{1}%

\section{La integral gaussiana}

\begin{teorema}[Integral gaussiana]\label{teo:gauss}
La integral impropia $\int_{-\infty}^{\infty}e^{-x^2}dx$ converge y vale $\sqrt\pi$.
\end{teorema}

\addtocounter{figure}{1}%

\begin{corolario}\label{cor:gauss-normal}
$\displaystyle\int_{-\infty}^{\infty}e^{-x^2/2}\,dx=\sqrt{2\pi}$. En consecuencia, la densidad normal estándar $\varphi(x)=\frac{1}{\sqrt{2\pi}}e^{-x^2/2}$ integra~$1$.
\end{corolario}

\addtocounter{definicion}{1}%

\addtocounter{definicion}{1}%

